\section{Step 35 Retain the outer byte walkers}\label{chacha-s35}

After successful preflight, the supplied 192-byte frame saves the outer working registers and return address. Keep source, destination and remaining count in \texttt{s0}, \texttt{s1} and \texttt{s2}; clear the state-copy offset. The frame's first 128 bytes hold a local state and keystream.

\subsection{Write and inspect}
Replace the body of region \texttt{stream\_setup}. Retain its \texttt{\# BEGIN} and \texttt{\# END} comments and the supplied labels outside them.

\Needspace{9\baselineskip}
\MARSCodeLabel{MIPS assembly}
\begin{minted}[samepage=true]{mars-idea}
    move $s0, $a1
    move $s1, $a2
    move $s2, $a3
    move $t0, $zero
\end{minted}

Use recipe \texttt{s35-1} with the cumulative program through this edit. Inspect its named checkpoint before continuing to the temporary or normal return.

\subsection{Expected observation}
The outer byte walkers and remaining count are retained; the local state copy is not yet initialized.

\Needspace{9\baselineskip}
\subsection{Check your understanding}
Why do the walkers need to survive a block call?

\begin{itemize}
  \item \textbf{A} The simulator automatically moves them to the heap.
  \item \textbf{B} The source must be word-aligned.
  \item \textbf{C} The nested block routine may change arguments and temporaries.
\end{itemize}
